\begin{scanpage}{0195}{177}
Here the quantity
\[
J_0=\frac{c\mathcal E}{4\ln\frac{i}{\gamma ka}},\quad\gamma=1.781\ldots\tag{28.04}
\]
determines the initial value of the current wave propagated from the excitation point$^{(1)}$. The function $\psi(z)$ is the solution of the integral equation, and in addition to the variable $z$ it also depends on the parameters $k$ and $a$. We will not list here all the properties of the function $\psi(z)$, but let us note only that it satisfies the conditions
\[
\psi(0)=1,\qquad\psi(\infty)=0,\tag{28.05}
\]
and its absolute value monotonically decreases with an increase of $z$. This decrease, which is rather slow and does not have an exponential character, is due to radiation.

The constants $A_1$ and $A_2$ determine the initial values of the current waves originating at the points $z=z_1$ and $z=z_2$, respectively, and travelling in the direction towards the opposite end of the conductor. These constants were found from the conditions at the conductor ends
\[
J(z_1)=J(z_2)=0\tag{28.06}
\]
and equal
\[
\left.\begin{aligned}A_1&=-\frac1D[\psi(-z_1)-\psi(z_2)\psi(L)e^{2ikz_2}]e^{-ikz_1},\\
A_2&=-\frac1D[\psi(z_2)-\psi(-z_1)\psi(L)e^{-2ikz_1}]e^{ikz_2},
\end{aligned}\right\}\tag{28.07}
\]
where
\[
D=1-\psi^2(L)e^{2ikL}.\tag{28.08}
\]
\par\bigskip\hrule\smallskip
$(1)$Footnote appears on page 216.
\end{scanpage}
